
% =====================================================================
\section{T0: the floor}
\label{sec:floor}
% =====================================================================

The first substructure is the floor itself: in a contact graph no part can be
separated from the rest at zero cost. The statement is immediate from the
premise, but its consequences---restated for cuts, for parts, and for the
limiting behaviour under refinement---are what the later rungs use, so we record
them with care.

\subsection{Every separation is costly}

\begin{theorem}[Floor Theorem]
\label{thm:floor}
Let $\Gph$ be a contact graph with floor $\floorc>0$. Then every nonempty proper
part $U\subsetneq\V$ has boundary weight bounded below by the floor:
\[
\boxed{\ \wt(\cut(U))\ \ge\ \floorc\ >\ 0.\ }
\]
In particular no part is separated from its complement by a cut of zero
weight: a \emph{sharp} separation, $\wt(\cut(U))=0$, is impossible.
\end{theorem}

\begin{proof}
Since $\Gph$ is connected (\Cref{ax:finite}) and $U$ is a nonempty proper
subset, there is at least one edge with one endpoint in $U$ and the other in
$\V\setminus U$; otherwise $U$ and $\V\setminus U$ would be unjoined, splitting
the connected graph $\Gph$ into two nonempty components, a contradiction. Hence
$\cut(U)\neq\varnothing$, and
\[
\wt(\cut(U))\;=\!\!\sum_{e\in\cut(U)}\!\!\wt(e)\;\ge\;\wt(e^\ast)\;\ge\;\floorc>0
\]
for any $e^\ast\in\cut(U)$, using $\wt(e)\ge\floorc$ from \Cref{ax:finite}. A
sharp separation would require $\wt(\cut(U))=0<\floorc$, impossible.
\end{proof}

\begin{definition}[Boundary cost; the realised floor]
\label{def:boundarycost}
The \emph{boundary cost} of a part $U$ is $b(U):=\wt(\cut(U))$. The
\emph{realised floor} of $\Gph$ is
\[
\floorc_\Gph\;:=\;\min_{\varnothing\neq U\subsetneq\V} b(U)\;\ge\;\floorc>0,
\]
the least boundary cost of any part. By \Cref{thm:floor} it is strictly
positive, and by finiteness of $\V$ the minimum is attained.
\end{definition}

\begin{corollary}[No zero-cost individuation]
\label{cor:no-zero}
In a contact graph, individuating any part from the rest deposits boundary cost
at least $\floorc_\Gph>0$. Equivalently, the map $U\mapsto b(U)$ is bounded
away from zero on nonempty proper parts.
\end{corollary}

\begin{proof}
Immediate from \Cref{thm:floor} and \Cref{def:boundarycost}.
\end{proof}

\subsection{Refinement lowers but never abolishes the floor}

A finer description splits parts into smaller ones. We show this can only lower
the realised floor toward, but never to, zero---so long as the graph remains
finite with a positive weight floor.

\begin{proposition}[Refinement bound]
\label{prop:refine}
Let $\Gph'$ be obtained from $\Gph$ by any sequence of vertex splits that
preserves the floor $\floorc$ on all edges \textup{(}each new edge again has
weight $\ge\floorc$\textup{)}. Then the realised floor satisfies
$\floorc\le\floorc_{\Gph'}$, and in particular $\floorc_{\Gph'}>0$. The realised
floor can decrease under refinement but is bounded below by $\floorc$ and never
reaches zero at any finite stage.
\end{proposition}

\begin{proof}
$\Gph'$ is again finite (a finite number of splits of a finite graph) and, by
hypothesis, has all edge weights $\ge\floorc$, hence is a contact graph with
floor $\floorc$. By \Cref{thm:floor} applied to $\Gph'$, every part of $\Gph'$
has boundary cost $\ge\floorc$, so $\floorc_{\Gph'}\ge\floorc>0$. Splitting can
only introduce parts with smaller boundary cost, so $\floorc_{\Gph'}\le
\floorc_\Gph$ is possible; in all cases the bound $\floorc_{\Gph'}\ge\floorc$
holds. Only in the non-admissible limit of infinitely many splits, or of edge
weights tending to $0$, would the bound vanish---both excluded by
\Cref{ax:finite}.
\end{proof}

\begin{remark}[What T0 does and does not assert]
\label{rem:floor-scope}
\Cref{thm:floor} does not assert that some ambient continuum has thick
boundaries; it asserts that a \emph{finite, positively weighted} graph cannot
exhibit a zero-cost cut. The floor is a property of the graph-with-its-weights,
i.e. of the finite system of separated parts, not of any underlying space. This
is the precise sense in which, for finitely many bounded parts, ``every
distinction costs something.'' Every later rung consumes \Cref{thm:floor} as the
guarantee that boundaries, residuals, and gated steps are bounded away from
zero.
\end{remark}
